\section{Exact collisions among inhomogeneous solutions}
\label{sec:exact-inhomogeneous}

The geodesic-image construction for Gowdy is established
\cite[\S\S III--IV]{berger}. Here we use it to compare what different
solutions reveal at the prescribed snapshots.
For an affinely parametrized target geodesic $\gamma$,
$U(t,\theta)=\gamma(w(t,\theta))$ solves \eqref{wave} if
\[
 w_{tt}+t^{-1}w_t-w_{\theta\theta}=0.
\]
If $\gamma$ has speed $\sigma$, the constraint sources for
$U$ are $-\sigma^2 t(w_t^2+w_\theta^2)$ and
$-2\sigma^2t w_t w_\theta$. In the examples below $w(a,\theta)$
is constant, so the periodic momentum constraint holds at $a$
and is propagated. Normalize the metric with one common
anchor at $a$.

For a Fourier mode $k\ge1$, define the solution with
$u_k(a)=0$ and $a u_k'(a)=1$ by
\begin{equation}
 u_k(t)=\frac{\pi}{2}
   [J_0(ka)Y_0(kt)-Y_0(ka)J_0(kt)].
 \label{eq:invisible}
\end{equation}
It has a terminal node when $u_k(b)=0$.
Every ratio $R=b/a>1$ admits such a pair of slices.
Indeed, take any positive eigenvalue $(ka)^2$ of the regular
Dirichlet problem $-(r f')'=(ka)^2 r f$ on $[1,R]$ and set
$a=(ka)/k$, $b=Ra$. Increasing the integer $k$ moves both slices
arbitrarily close to $t=0$ at fixed ratio.
The logarithmic derivative at the singular limit satisfies
$\lim_{t\downarrow0}t u_k'(t)=J_0(ka)$.

\begin{theorem}
\label{th:exact-inhomogeneous}
For every ratio $R>1$ there are finite snapshots $a,b$ with
$b/a=R$ and distinct smooth inhomogeneous vacuum Gowdy solutions
having identical induced spatial metrics and identical $Q_t$ at
both times. Their metric potentials $\lambda$ agree at every
$(t,\theta)$, while their $P_t(a)$ differ. These pairs occur
arbitrarily close to a constant wave map on the observation
window and have uniformly subcritical asymptotic speed.
\end{theorem}

\begin{proof}
Let $(P_0,Q_0)$ be a common initial point. Set
$w=\epsilon u_k(t)\cos(k\theta)$ and choose the two geodesics
with initial tangent components $(g,q)$ and $(-g,q)$ in the
orthonormal frame $(\partial_P,e^{-P_0}\partial_Q)$, with
$g,q\ne0$. The geodesics have the same speed
$\sigma=(g^2+q^2)^{1/2}$. At the two nodes the fields $P,Q$
return to $(P_0,Q_0)$, and in each solution
$Q_t=e^{-P_0}q w_t$ there. Their initial $P_t$ are
$\pm g\epsilon\cos(k\theta)/a$. The two sources of $\lambda$
are identical at every time by equal speeds and the shared
scalar phase. The anchors therefore give the same $\lambda$.
The solutions have nonconstant orbit fields at generic interior
times. As $\epsilon\to 0$ the data and their regular-time
evolutions approach the constant background in every fixed
Sobolev norm. Their asymptotic speeds are
$\sigma|\epsilon J_0(ka)\cos(k\theta)|$, uniformly below one
for small enough $\epsilon$.
\end{proof}

Any finite collection of snapshots selected from the nodes
of this one mode has the same ambiguity. At a fixed $a$,
the resonant terminal times are countable and dense in $(a,\infty)$.
For density, rescale $u_k$ by $\sqrt t$ and apply the
Volterra formula on a compact positive time interval.
After normalizing the initial derivative to $k$, the
transfer is $\sin(k(t-a))+O(k^{-1})$, uniformly on the
interval. Opposite signs on either side of nearby free nodes
give exact zeros by continuity.

\begin{proposition}
\label{prop:third}
Restrict to geodesic-image solutions with $U(a,\theta)=U_0$
spatially constant and with a common initial metric anchor.
Full orbit fields $P,Q$ at later times $t_1,\dots,t_m$
determine the solution if and only if no positive Fourier
mode satisfies $u_k(t_1)=\cdots=u_k(t_m)=0$.
In particular, at fixed $a,b$ a third full snapshot at any
$c\in(a,b)$ outside a countable set resolves the ambiguity
in this class.
\end{proposition}

\begin{proof}
If every measured field equals $U_0$, a unit-speed signed
arclength parameter along any candidate geodesic has
constant zero data at each measured slice and at $a$.
Its zero Fourier mode evolves from zero by a multiple of
$\log(t/a)$ and hence vanishes. Each other mode is a multiple
of $u_k(t)$, so each vanishes if there is no common node.
If there is a measured point different from $U_0$, the unique
complete target geodesic through those two points becomes
known. Choose one orientation and read the signed arclength
at all snapshots. The same mode argument gives uniqueness.
The constraints then recover $\lambda$ with the anchor.
A common node gives the pair in Theorem~\ref{th:exact-inhomogeneous}.
For any $k$ with $u_k(b)=0$, its additional zero times
are discrete in $(a,b)$. Their countable union gives the
exceptional set.
\end{proof}

\begin{proposition}
\label{prop:safe-gap}
For each $a>0$ and each $b>a$, a full orbit snapshot
at a specified $c\in(a,b)$ gives a quantitative
inverse on the constant-initial-position geodesic-image
class. Set $M=1/(4a^2)$ and choose an integer
$L>\max\{\pi/(b-a),\,2M+2\}$.
Take $c=a+\pi/L$.
Using just the full orbit snapshots at $a$ and $c$,
recovery of the initial velocity in energy normalization
has sharp loss one. In particular, for each real $s$,
\[
 \|U_t(a)-\widehat U_t(a)\|_{H^{s-1}}
 \le C(a,c)\|
       \log_{U_0}U(c)-\log_{U_0}\widehat U(c)
      \|_{H^{s+1}} .
\]
Both solutions start at the same constant target point $U_0$.
The velocity difference is taken in $T_{U_0}\mathbb H^2$.
\end{proposition}

\begin{proof}
The Riemannian logarithm turns a geodesic through $U_0$
into a straight line in $T_{U_0}\mathbb H^2$;
therefore $Z=\log_{U_0}U$ satisfies the scalar
equation componentwise, also when comparing two
different geodesics. On frequency $k\ge1$, rescale
by $\sqrt t$ and normalize the initial slope to $k$.
Writing the resulting solution as
$z=\rho\sin\varphi$, $z'=k\rho\cos\varphi$ gives
\[
 \varphi'=k+\frac{W(t)}{k}\sin^2\varphi,\qquad
 \frac{\rho'}{\rho}=-\frac{W(t)}{2k}\sin(2\varphi),
 \qquad W(t)=\frac1{4t^2}.
\]
The starting values are $\rho(a)=1$, $\varphi(a)=0$.
With $d=\pi/L$, the correction to the free phase at $c$
is strictly between zero and $Md/k$. If $k>M$,
this is less than the spacing $d$ of the free phases
modulo $\pi$. No mode can reach a zero. If $k\le M$,
the phase is between zero and $(k+M/k)d\le 2Md<\pi$.
There are thus no finite-frequency zeros.

For large $k$ the same phase bounds hold at every time
in the window, so
\[
 \varphi(c)-kd
 =\frac1k\int_a^c W(t)\sin^2(k(t-a))\,dt+O(k^{-2})
 =\frac1{2k}\int_a^c W(t)\,dt+O(k^{-2}).
\]
The second equality uses integration by parts.
The amplitude is $1+O(k^{-1})$.
On multiples of $L$, the multiplier is therefore of
order $k^{-1}$ with a nonzero leading coefficient.
On large nonmultiples it is bounded away from zero.
The zero-mode transfer is $a\log(c/a)$
for the unscaled velocity. This proves the estimate.
Taking single high modes on one target geodesic shows
sharpness for the pair $a,c$.
\end{proof}


The same construction also extends the homogeneous fold
to histories with strictly positive asymptotic speed.
For the second row of Table~\ref{tab:collisions}, take
$w=\log(t/a)+\epsilon u_k(t)\cos(k\theta)$
with the same Dirichlet placement and use the two colliding
geodesics from that row. The common initial $Q_t$ is
$e^{-P(a)}q[1+\epsilon\cos(k\theta)]/a$.
The phases at the two snapshots are still $0$ and $\ell$,
so the four original fields agree. The starting $\lambda$
anchor agrees as well. The differing initial $P_t$ guarantee
distinct inhomogeneous histories. The asymptotic speeds
$\sigma_j|1+\epsilon J_0(ka)\cos(k\theta)|$ are uniformly
in $(0,1)$ for small enough $\epsilon$. These lift the
partial-record fold; the two terminal $Q$ values differ.

In an illustrative full-metric collision with $b/a=e$,
$k=4$, $\epsilon=0.02$, $g=0.55$ and $q=0.37$,
the reference time is $a\simeq0.451633$.
The two sets of snapshot data agreed to floating-point
precision and the initial $P_t$ differed by $0.04871$
in supremum norm. Reading $P$ at the midpoint separated
the two by $0.00876$. In a separate check of the partial-record
lift at $b/a=e^3$, the two computed asymptotic-speed ranges
were $[0.9311,0.9689]$ and $[0.6530,0.6796]$.
Both agreements follow from the exact formulas.
